% Generated draft: requires mathematical and source review.
\documentclass[11pt]{article}
\usepackage[margin=25mm]{geometry}
\usepackage{fontspec}
\setmainfont{Latin Modern Roman}
\usepackage{amsmath,amssymb,mathtools,unicode-math}
\setmathfont{Latin Modern Math}
\usepackage{enumitem,needspace,xurl,hyperref}
\setlength{\parindent}{0pt}
\setlength{\parskip}{5pt}
\setlength{\emergencystretch}{3em}
\newcommand{\N}{\mathbb{N}}
\newcommand{\Z}{\mathbb{Z}}
\newcommand{\Q}{\mathbb{Q}}
\newcommand{\R}{\mathbb{R}}
\newcommand{\C}{\mathbb{C}}
\newcommand{\F}{\mathbb{F}}
\newcommand{\A}{\mathbb{A}}
\newcommand{\PP}{\mathbb{P}}
\newcommand{\E}{\mathbb{E}}
\newcommand{\eps}{\varepsilon}
\newcommand{\dd}{\,\mathrm{d}}
\begin{document}
\section*{UnsolvedMath ID: 30004477}
\textbf{Source problem code:} catalogue entry
\subsection{Definitions and mathematical statement}
Fix embeddings \(K\subset\overline{\Q}\subset\C\), where \(K\) is a number field. A polarized variation of integral Hodge structure consists of a locally constant sheaf \(\mathbb V_{\Z}\) of finite-rank free abelian groups, a Hodge filtration on the associated holomorphic vector bundle, and a flat polarization. Its fibers are polarized pure Hodge structures of a fixed weight, and its filtration satisfies Griffiths transversality
\[
\nabla F^p\subseteq F^{p-1}\otimes\Omega^1.
\]
The record does not specify the weight, the hypotheses on the base, or the precise descent data meant by saying that this variation is defined over \(K\).

Under the usual Mumford--Tate interpretation, a special subvariety is an irreducible closed algebraic subvariety maximal for its generic Mumford--Tate group, with these groups compared using local parallel transport. The Mumford--Tate group of a rational Hodge structure is the smallest algebraic subgroup over \(\Q\) of its general linear group through which its defining real Hodge homomorphism factors. This interpretation of special subvariety is provisional: the supplied record does not define the term or specify its conventions.

\textbf{Supplied conjecture, with the descent conclusion made explicit.} Let \(\mathbb V\) be a polarized variation of integral Hodge structure defined over a number field \(K\). Write \(S_K\) for its intended algebraic base over \(K\), and put \(S_{\C}=S_K\times_K\C\). For every associated special subvariety \(Z\subseteq S_{\C}\), there should exist a finite extension \(K\subseteq L\subseteq\overline{\Q}\) and a closed subvariety \(Z_L\subseteq S_K\times_K L\) such that
\[
Z=Z_L\times_L\C.
\]
Moreover, writing \(Z_{\overline{\Q}}=Z_L\times_L\overline{\Q}\), one should have
\[
\text{for every }\sigma\in\operatorname{Gal}(\overline{\Q}/K),\qquad
\bigl(\sigma Z_{\overline{\Q}}\bigr)\times_{\overline{\Q}}\C
\text{ is special for }\mathbb V\text{ on }S_{\C}.
\]
Here conjugation acts on the algebraic coefficients inside the fixed base defined over \(K\). The orbit consists of finitely many distinct conjugate subvarieties because a field of definition \(L\) is finite over \(K\); the Galois group itself need not be finite.

The notation \(S_K\) exposes base-descent data needed for this formulation; it does not supply the missing hypotheses or a definition of descent for the variation. Those points, and the intended meaning of special subvariety, require confirmation from the cited report. The record presents this assertion as a conjecture with supplied status open.
\subsection{Short English statement}
The supplied conjecture asks whether every special subvariety associated with a polarized integral variation of Hodge structure defined over a number field \(K\) descends to a finite extension of \(K\), and whether each of its finitely many distinct Galois conjugates is again special for the same variation. The exact base hypotheses and the meanings of descent and special subvariety require clarification.
\subsection{Sources}
\begin{itemize}
\item Local UnsolvedMath registry, permanent ID \(30004477\), title \emph{Galois Descent of Special Hodge Subvarieties}; exact source problem code \texttt{catalogue entry}. The supplied collection is \emph{Oberwolfach Reports Open Problems}; its proposed-year field is \(2020\). No individual proposer is supplied.
\item Supplied report citation: \emph{Arithmetic Geometry (2021)}, Oberwolfach Reports, DOI \texttt{10.4171/owr/2020/20}, pp.\ 1023--1082. Source URL and supplied literature-source link, labeled ``original OWR report'': \url{https://doi.org/10.4171/owr/2020/20}. The background describes Oberwolfach Reports as documenting mathematical workshops and their open-problem discussions.
\item Supplied literature assessment dated \texttt{2026-08-22}: the general Galois-descent assertion remains conjectural, while work attributed to Voisin and Saito--Schnell provides sufficient criteria rather than a proof of the full assertion. This is supplied dated curation derived from the OWR record, not an independently verified literature assessment. No specific references for those authors' results are supplied.
\item The registry labels the statement \texttt{corrected\_verified} and reports that it was checked against the official Oberwolfach report. This is a supplied verification claim; the report and that claim have not been independently checked for this draft.
\end{itemize}
\end{document}
